\documentclass[11pt]{article}

\usepackage[margin=1in]{geometry}
\usepackage{amsmath,amssymb,amsthm}
\usepackage{enumitem}
\usepackage{hyperref}
\usepackage{listings}
\usepackage[T1]{fontenc}
\usepackage[scaled]{beramono}
\usepackage{courier}
\usepackage{microtype}

\lstdefinestyle{code}{
  basicstyle=\ttfamily\small,
  breaklines=true,
  frame=single,
  columns=fullflexible,
  keepspaces=true,
  showstringspaces=false
}

\title{%
  L0 Scope-Bound Non-Repudiation \\ and Tamper-Evident Audit Trails\\
  in Phase Mirror ACE and Witness Protocols\\[0.5ex]
  \large Defensive Publication and Design Report
}
\author{Citizen Gardens \\ The Foundation of Multiplicity}
\date{April 25, 2026}

\begin{document}
\maketitle

\begin{abstract}
This report describes two tightly coupled design developments in the
Phase Mirror ecosystem:

\begin{enumerate}[label=(\alph*)]
  \item a Level~0 (\emph{L0}) scope-binding invariant for Ed25519 witness
        signatures, including a stratified non-repudiation model and a
        two-root trust architecture based on SPIFFE SVIDs; and
  \item a tamper-evident, testable audit trail for the Arithmetic Control
        Engine (ACE) in the PIRTM project, including HMAC chaining,
        deterministic testing under clock injection, float-normalization
        gates, and thread-safety.
\end{enumerate}

The intent is explicitly defensive: to place in the public domain a
precise combination of invariants, message formats, and code patterns
that enable provable non-repudiation and tamper-evidence in a
governance-grade system. The text is written as a comprehensive
prior-art publication and as a design reference for implementers.
\end{abstract}
\newpage
\tableofcontents
\newpage
\section{Executive Summary}

This report documents a sequence of design decisions in two domains:

\begin{itemize}
  \item \textbf{Scope-bound signatures (L0-SIG-SCOPE).} Ed25519 witness
        signatures are required to be cryptographically bound to a
        well-defined \emph{scope}: domain, intent class, audience
        (both class and instance), governance surface (tenant/jurisdiction),
        time window, replay protection, and key purpose. Signatures that
        lack these bindings are rejected at the verifier. The scope
        enforcement is stratified by a non-repudiation level (\texttt{NRLevel}),
        and executor identity is anchored in a platform attestation
        trust root (SPIFFE SVIDs).

  \item \textbf{Tamper-evident audit trails (ADR-012 for ACE).}
        The ACE subsystem maintains an append-only, HMAC-chained log
        of certificates, with the HMAC key derived from the previous
        chain hash. The design eliminates non-determinism sources
        (clock and session identifier) via constructor injection,
        normalizes float fields at the serialization boundary, rejects
        non-finite values loudly, and achieves thread-safe atomic
        updates of the chain head.
\end{itemize}

Both threads share the same central tension: the boundary between
\emph{validity} (what actions and values are acceptable in the first
place) and \emph{defense in depth} (how failures are surfaced and
made tamper-evident). The key principle is that upstream gate failures
(e.g.\ loss of spectral gap, missing scope bindings) must be detected
at the earliest possible layer, and that downstream cryptographic
mechanisms (HMAC, signatures, anchors) must never silently encode
invalid state.

\section{Mathematical and Cryptographic Overview}

\subsection{HMAC-Chained Audit Trail}

The ACE audit trail is an append-only log of certificate entries
\(\{e_i\}_{i=1}^n\) with the following fields:
\[
  e_i = \bigl(\text{prime\_index}_i,\ \text{cert}_i,\ \text{wall\_clock}_i,\
         \text{session\_id},\ \text{chain\_hash}_i\bigr).
\]

Let \(\mathcal{H}\) denote HMAC-SHA256. The chain is defined by:
\begin{align}
  k_0 &= \underbrace{\text{bytes.fromhex}(\text{``0''}\times 64)}_{\text{32 zero bytes}},\\
  k_i &= \text{bytes.fromhex}(\text{chain\_hash}_{i-1}), \quad i \ge 1,\\
  m_i &= \text{JSON\_CANONICAL}\left(
           \text{prime\_index}_i,\ \text{cert}_i,\ \text{wall\_clock}_i,\ \text{session\_id}
        \right), \\
  d_i &= \mathcal{H}(k_i, m_i), \\
  \text{chain\_hash}_i &= \text{hex}(d_i).
\end{align}

Here \(\text{JSON\_CANONICAL}\) denotes UTF-8 JSON encoding with
sorted keys and \texttt{allow\_nan=False}, so that non-finite floats
(\(\pm\infty, \text{NaN}\)) are rejected with a \texttt{ValueError}
at the encoder boundary rather than silently encoded.\footnote{See
Python's \texttt{json.dumps} and \texttt{allow\_nan} semantics
\cite{web:32,web:34}.}

Verification replays the chain from the fixed genesis key \(k_0\).
Any modification to a logged entry \(e_j\) changes the recomputed
digest \(d_j\), causing a mismatch at or before index \(j\).

\subsection{Scope-Bound Signatures}

We model the signed message \(M\) for an Ed25519 witness signature as
a concatenation of length-prefixed byte strings:
\[
  M = \bigl\|_{i=1}^k \ell_i \parallel x_i,
\]
where each field \(x_i\) is a UTF-8 string, and \(\ell_i\) is a fixed
size length prefix (e.g.\ 4-byte big-endian). This avoids ambiguity
in boundary parsing.

The fields in order are:

\begin{enumerate}[label=(\roman*)]
  \item Domain separator with IntentClass Registry (ICR) pin:
        \(\texttt{"PHASE\_MIRROR/WITNESS/v1/icr:}\langle \texttt{icr\_version}\rangle\texttt{"}\).
  \item \(\texttt{intentClass}\): a closed enum from the ICR.
  \item \(\texttt{scopeHash} = \text{SHA256}(\text{RFC8785}(\text{scope\_object}))\),
        where the scope object includes tenant, project, jurisdiction,
        consent, and policy version.
  \item \(\texttt{aud.class}\) and \(\texttt{aud.version}\), always present.
  \item \(\texttt{aud.instance}\): present and non-empty if and only if
        the non-repudiation level \( \texttt{NRLevel} = \texttt{NR\_FULL} \).
  \item \(\texttt{intentHash} = \text{SHA256}(\text{RFC8785}(\text{cap\_object}))\).
  \item \(\texttt{traceHash}\): hash of the trace atom being witnessed.
  \item Temporal and anti-replay fields: \(\texttt{iat}, \texttt{exp}, \texttt{nonce}\).
  \item \(\texttt{kid}\): key identifier carrying key purpose.
\end{enumerate}

The witness signature is then:
\[
  \sigma = \text{Ed25519.Sign}(sk, M),
\]
and the verifier reconstructs \(M\) and checks
\(\text{Ed25519.Verify}(pk, M, \sigma)\) after enforcing the structural
and trust-root invariants described in Section~\ref{sec:l0-scope}.

\subsection{Stratified Non-Repudiation via NRLevel}

We introduce an enum \(\texttt{NRLevel} \in
\{\texttt{NR\_FULL},\ \texttt{NR\_CLASS},\ \texttt{NR\_NONE}\}\)
assigned per capability action (\(\texttt{cap.action}\)) by a versioned,
governed registry (the ICR). This field determines which non-repudiation
invariants are required.

\begin{itemize}
  \item \(\texttt{NR\_FULL}\): governance-relevant actuation; requires
        SPIFFE SVID-based executor identity, instance-level audience
        binding, and historical SVID key archival.
  \item \(\texttt{NR\_CLASS}\): non-governance actuation; class-level
        binding suffices; instance binding is not enforced.
  \item \(\texttt{NR\_NONE}\): witness, receipt, or trace-only operations;
        no instance binding, evidentiary only.
\end{itemize}

The verifier enforces different invariants depending on \(\texttt{NRLevel}\);
see Section~\ref{sec:l0-scope}.

\section{ACE Audit Trail Design (ADR-012)}

\subsection{Design Invariants}

The ACE audit trail implementation, \texttt{AceAuditTrail}, enforces
the following invariants:

\begin{enumerate}[label=(I\arabic*)]
  \item \textbf{NO\_CONTEXT\_MUTATION:} entries are never modified after
        \texttt{append()}.
  \item \textbf{DETERMINISTIC\_UNDER\_INJECTION:} all non-determinism
        sources relevant to the HMAC chain (wall-clock, session ID)
        are injectable at construction so that tests can compute exact
        chain hashes via equality.
  \item \textbf{TAMPER\_EVIDENT:} \texttt{verify\_chain()} replays the
        HMAC chain from the genesis hash and returns \texttt{False} on
        any single-field modification.
  \item \textbf{HKDF\_COMPOSABLE:} the HMAC key at each step is
        \texttt{bytes.fromhex(\_prev\_hash)}, i.e.\ raw 32-byte digest
        material with full byte-value range \([0x00,\dots,0xFF]\),
        zero-padded by the HMAC block schedule. Hex encoding is used
        only for storage and display.
  \item \textbf{THREAD-SAFE:} concurrent \texttt{append()} and
        \texttt{verify\_chain()} calls are serialized around the
        \((\_entries, \_prev\_hash)\) mutation pair with a lock;
        verification uses a snapshot.
\end{enumerate}

\subsection{Clock Injection and Determinism}

In naive implementations, \texttt{datetime.now(timezone.utc)} is called
inside \texttt{append()}, making both the wall-clock and the derived
\texttt{chain\_hash} non-deterministic. This prevents exact equality
tests on the HMAC chain.

The corrected design injects a clock callable at construction:

\begin{lstlisting}[style=code,language=Python]
class AceAuditTrail:
    def __init__(
        self,
        session_id: Optional[str] = None,
        *,
        clock: Optional[Callable[[], datetime]] = None,
    ) -> None:
        self.session_id = session_id or str(uuid.uuid4())
        self._clock = clock or (lambda: datetime.now(timezone.utc))
        self._entries: list[AuditEntry] = []
        self._prev_hash: str = _GENESIS_HASH
        self._lock = threading.Lock()
\end{lstlisting}

Tests construct the trail with a fixed clock and session ID, enabling
exact \texttt{==} assertions on \texttt{chain\_hash}.

\subsection{HMAC Key Material and One-Shot Digest}

The HMAC key is derived from the previous chain hash via
\texttt{bytes.fromhex()}, not via UTF-8 encoding of a hex string:

\begin{lstlisting}[style=code,language=Python]
chain_hash = hmac.digest(
    bytes.fromhex(self._prev_hash),   # 32 raw bytes
    entry_json.encode(),
    hashlib.sha256,
).hex()
\end{lstlisting}

Using \texttt{bytes.fromhex()} ensures the key bytes span the full
range and interact with HMAC's block-size reduction as intended,
preserving compatibility with HKDF-Extract when the trail tail is
used as input key material \cite{web:16,web:22}.

\subsection{Float Normalization and Non-Finite Values}

\texttt{AceCertificate} carries several float fields
(lipschitz bounds, gap lower bounds, tail bounds). To ensure JSON
encoding is both deterministic and standards-compliant, the
serialization layer enforces finiteness:

\begin{lstlisting}[style=code,language=Python]
class AuditSerializationError(ValueError):
    ...

def _assert_finite(value: float, field: str) -> float:
    if not math.isfinite(value):
        raise AuditSerializationError(
            f"certificate_to_dict(): field '{field}' is non-finite "
            f"({value!r}). ... Root cause: upstream certifier gate "
            f"condition did not return certified=False when gap_lb < delta."
        )
    return value

_FLOAT_FIELDS = (..., "tail_bound", ...)

def certificate_to_dict(cert: AceCertificate) -> dict:
    d = dataclasses.asdict(cert)
    if "level" in d and isinstance(d["level"], CertLevel):
        d["level"] = d["level"].value
    for field in _FLOAT_FIELDS:
        if field in d and d[field] is not None:
            d[field] = _assert_finite(d[field], field)
    return d
\end{lstlisting}

At the audit trail boundary, \texttt{json.dumps} is called with
\texttt{allow\_nan=False}, ensuring any non-finite float that escaped
\texttt{\_assert\_finite} (e.g.\ inside a nested \texttt{details}
dictionary) triggers a \texttt{ValueError}:

\begin{lstlisting}[style=code,language=Python]
entry_json = json.dumps(entry_data, sort_keys=True, allow_nan=False)
\end{lstlisting}

Upstream certifier functions (\texttt{certify\_l3}, \texttt{certify\_l4},
\texttt{certify\_l5}) are responsible for returning
\texttt{certified=False} when gap conditions fail, and the
\texttt{AceProtocol} layer omits uncertified certificates from the
trail entirely.

\subsection{Thread Safety and Immutability}

The audit trail uses a \texttt{threading.Lock} to serialize the
critical region:

\begin{lstlisting}[style=code,language=Python]
def append(self, cert: AceCertificate, prime_index: int) -> AuditEntry:
    ...
    entry_json = json.dumps(entry_data, sort_keys=True, allow_nan=False)

    with self._lock:
        chain_hash = hmac.digest(
            bytes.fromhex(self._prev_hash),
            entry_json.encode(),
            hashlib.sha256,
        ).hex()
        entry = AuditEntry(..., chain_hash=chain_hash)
        self._entries.append(entry)
        self._prev_hash = chain_hash
    return entry
\end{lstlisting}

Verification obtains a snapshot of entries under the lock, then
replays HMACs without holding the lock:

\begin{lstlisting}[style=code,language=Python]
def verify_chain(self) -> bool:
    with self._lock:
        snapshot = list(self._entries)
    prev = _GENESIS_HASH
    for e in snapshot:
        entry_data = {...}
        entry_json = json.dumps(entry_data, sort_keys=True, allow_nan=False)
        expected = hmac.digest(bytes.fromhex(prev), entry_json.encode(),
                               hashlib.sha256).hex()
        if not hmac.compare_digest(expected, e.chain_hash):
            return False
        prev = e.chain_hash
    return True
\end{lstlisting}

The \texttt{AuditEntry} dataclass is made \emph{frozen} to prevent
accidental mutation in production, but tamper tests deliberately
bypass this via \texttt{object.\_\_setattr\_\_} to demonstrate that
the integrity boundary is the HMAC chain, not Python's object model.

\section{L0 Scope-Bound Signatures and Executor Identity}
\label{sec:l0-scope}

\subsection{Decomposing Audience and Intent}

The initial invariant treated \texttt{aud} as a single policy dial.
This is insufficient. The correct abstraction decomposes audience
into:
\begin{itemize}
  \item \texttt{aud.class}: logical adapter type (e.g.\ text executor).
  \item \texttt{aud.instance}: concrete executor instance identity.
  \item \texttt{aud.version}: adapter API contract version.
\end{itemize}

At the same time, the \emph{intent} being authorized is classified
into an \texttt{intentClass} drawn from a versioned registry:

\begin{itemize}
  \item \texttt{actuation}: state-changing operation.
  \item \texttt{witness}: observation only.
  \item \texttt{receipt}: acknowledgement only.
  \item \texttt{trace-atom}: pure telemetry/provenance.
\end{itemize}

The key insight is that the requirement for instance-level binding
depends on \texttt{intentClass}: it is mandatory for actuation but
not for witness-only operations.

\subsection{IntentClass Registry and Non-Repudiation Level}

The IntentClass Registry (ICR) is a JSON artifact mapping
\(\texttt{cap.action}\) to \(\texttt{intentClass}\) and \(\texttt{NRLevel}\).
Crucially, the registry is versioned and its version is embedded in
the domain separator of the signed message.

Each action is assigned one of:
\begin{itemize}
  \item \texttt{NR\_FULL}: instance-level non-repudiation with
        SPIFFE SVID required.
  \item \texttt{NR\_CLASS}: class-level non-repudiation; instance
        binding optional.
  \item \texttt{NR\_NONE}: evidentiary-only.
\end{itemize}

The ICR also includes a governed field
\texttt{unclassified\_action\_default} with allowed values
\{\texttt{REJECT}, \texttt{NR\_CLASS}, \texttt{NR\_FULL}\},
which determines the fail-safe behavior when a \texttt{cap.action}
has no explicit entry. The default shipped value is \texttt{REJECT}.

\subsection{Sub-Invariants L0-SIG-SCOPE-*}

The original L0-SIG-SCOPE invariant unfolds into the following
sub-invariants:

\begin{description}
  \item[L0-SIG-SCOPE-CLASS] applies to all signatures. The verifier
    rejects any token with missing or mismatched \texttt{aud.class}
    or \texttt{aud.version}.

  \item[L0-SIG-SCOPE-INSTANCE] applies when
    \(\texttt{intentClass} = \texttt{actuation}\) and
    \(\texttt{NRLevel} = \texttt{NR\_FULL}\). The executor (or a
    gate on its behalf) rejects tokens if \texttt{aud.instance}
    is absent or does not match the executor's SPIFFE ID.

  \item[L0-SIG-SCOPE-INTENTCLASS] applies to all signatures. The
    verifier uses the ICR to determine the canonical
    \(\texttt{intentClass}\) for \(\texttt{cap.action}\), and rejects
    tokens where the signer-declared \texttt{intentClass} differs.
    This closes the classification oracle attack in which a signer
    might mislabel an actuation as a witness to evade stricter
    checks.

  \item[L0-SIG-SCOPE-SVID] applies when
    \(\texttt{NRLevel} = \texttt{NR\_FULL}\). The executor's SVID
    must validate against the SPIRE trust bundle at the time of
    execution (using the certificate valid at \texttt{iat}), and
    the key purpose \texttt{kid.type} must match
    \(\texttt{intentClass}\).
\end{description}

\subsection{Two-Root Trust Architecture via SPIFFE SVIDs}

Non-repudiation is anchored in two independent trust roots:

\begin{itemize}
  \item \textbf{Principal trust root:} a DID-based key hierarchy
        used to sign class-scoped delegation tokens. This root governs
        which actions are authorized and for which adapter classes.

  \item \textbf{Executor trust root:} a platform attestation authority
        (e.g.\ SPIRE) which issues short-lived SVIDs binding a
        SPIFFE ID to a workload instance based on runtime attestation
        (node, workload, TEE, etc.) \cite{web:77,web:90,web:104}.
\end{itemize}

The principal cannot safely assert instance identity, because it has
no direct observation of runtime placement. The executor cannot
self-assert instance identity without circularity. SPIRE breaks this
impasse by issuing identities to workloads based on attested runtime
properties, independent of the principal and the workload code.

Execution receipts bind both roots together by including:

\begin{itemize}
  \item the executor's SPIFFE ID and SVID key identifier;
  \item a hash of the delegation token;
  \item the non-repudiation level; and
  \item the Ed25519 signature over the scope-bound message.
\end{itemize}

\subsection{Verifier Algorithm}

The verifier executes the checks in the following order:

\begin{enumerate}[label=Step~\arabic*:]
  \item \textbf{Intent class and ICR version.} Resolve
        \(\texttt{intentClass}\) from the ICR using \(\texttt{cap.action}\).
        Reject if the token's declared \(\texttt{intentClass}\) does
        not match, or if the domain separator's ICR version differs
        from the registry.

  \item \textbf{Audience class and version.} Reject if
        \(\texttt{aud.class}\) or \(\texttt{aud.version}\) are absent
        or not in the verifier's supported set.

  \item \textbf{Instance binding (NR\_FULL).} If the action's
        \(\texttt{NRLevel} = \texttt{NR\_FULL}\), reject if
        \(\texttt{aud.instance}\) is absent or does not equal the
        executor's SPIFFE ID.

  \item \textbf{SVID validation (NR\_FULL).} Validate the executor's
        SVID against the SPIRE trust bundle and historical key archive
        using \(\texttt{svid\_key\_id}\) and \(\texttt{iat}\).
        Reject if expired at execution time, if issued by a different
        trust domain, or if \(\texttt{kid.type}\) does not match
        the ICR-specified purpose for the \(\texttt{intentClass}\).

  \item \textbf{Signature and replay.} Reconstruct the signed message
        \(M\); verify the Ed25519 signature with the appropriate key;
        reject on signature failure, replay (seen \texttt{nonce}), or
        invalid time window (\(\texttt{exp} < \texttt{now}\) or
        \(\texttt{iat}\) too far in the future).
\end{enumerate}

\subsection{ICR Governance and Empty Registry Semantics}

The IntentClass Registry includes a governance field
\texttt{unclassified\_action\_default} controlling how unlisted
\(\texttt{cap.action}\) values are treated:

\begin{itemize}
  \item \texttt{REJECT}: fail-safe default. Any unlisted action is
        rejected; CI gates fail if the registry has no classified
        actions and the default is \texttt{REJECT}.

  \item \texttt{NR\_CLASS}: explicit declaration that unlisted actions
        are treated as \texttt{NR\_CLASS}. An empty \texttt{actions[]}
        with this default describes an intentional NR\_CLASS-only
        deployment, which CI gates may allow.

  \item \texttt{NR\_FULL}: strict default; any unlisted action is
        treated as requiring SVID and instance binding.
\end{itemize}

CI gates read this field to distinguish between a truly unclassified
registry and an intentional NR\_CLASS-only deployment, preventing
the ``empty registry passes by accident'' failure mode.

\section{Defensive Publication Impact}

The combination of patterns described here is the subject of this
defensive publication:

\begin{itemize}
  \item A specific HMAC-chained audit trail that:
    \begin{itemize}
      \item injects time and session non-determinism at construction,
            not at append time;
      \item uses \texttt{bytes.fromhex} of prior chain hash as the
            HMAC key and \texttt{hmac.digest} as the one-shot form;
      \item enforces float finiteness and rejects non-finite values
            before JSON encoding;
      \item uses \texttt{allow\_nan=False} as a defense-of-last-resort;
      \item achieves thread-safe atomicity of log append and chain head
            update using a narrow lock.
    \end{itemize}
  \item A scope-bound Ed25519 signature scheme that:
    \begin{itemize}
      \item encodes domain, ICR version, intent class, audience
            class, audience instance, scope hash, intent hash, trace
            hash, time window, nonce, and key ID;
      \item uses a closed, versioned intent-class registry with
            a non-repudiation level per action;
      \item stratifies SPIFFE SVID requirements based on non-repudiation
            level;
      \item prevents classification oracle attacks by cross-checking
            declared intent class against registry-resolved class; and
      \item ties executor identity to a platform attestation authority
            rather than to either the principal or the workload.
    \end{itemize}
\end{itemize}

These elements, taken together, constitute prior art for any system
that claims an identical combination of invariants, message fields,
and code-level patterns in a governance-grade context.

\section{Recommendations and Future Work}

\begin{itemize}
  \item \textbf{CBOR migration for audit trail.} When cross-runtime
        verifiers (e.g.\ Go or Rust) are introduced, the audit trail
        should migrate from JSON to canonical CBOR (CTAP2 profile) for
        float encoding and cross-language determinism.

  \item \textbf{Air-gapped SPIRE deployments.} For environments without
        internet access, deploy SPIRE in an air-gapped configuration
        and document SVID trust-bundle distribution and archival.

  \item \textbf{Formal verification hooks.} Both the audit trail and
        the signature scope invariants lend themselves to formalization
        (e.g.\ in TLA+ or Coq), particularly the replay and
        non-repudiation properties.

  \item \textbf{Plane boundary enforcement.} The separation between
        Plane A (compute) and Plane C (FP calibration) identity models
        should be monitored as new features are added, ensuring that
        organizational identity (Plane C) is never conflated with
        executor identity (Plane A).
\end{itemize}

\appendix
\section{Mathematical Appendix}
\label{app:math}

This appendix collects the explicit proofs and norm bounds that were
referenced in the main text but omitted for brevity.

\subsection{Lyapunov decrease condition (L5)}
\begin{theorem}[Asymptotic stability from Lyapunov decrease]\label{thm:lyap}
Let a discrete-time system be described by
\[
  x_{k+1}=K\,x_k,
\]
with state $x_k\in\mathbb{R}^n$ and constant matrix $K$.
Suppose there exists a symmetric positive-definite matrix $P\succ0$
such that
\[
  \Delta V(x)=x_k^\top\!\bigl(K^\top P K-P\bigr)\,x_k<0
\quad\text{for all }x_k\neq0.
\tag{1}
\]
Then the origin $x=0$ is asymptotically stable.
\end{theorem}
\begin{proof}
Define the quadratic Lyapunov function $V(x)=x^\top P x$.
Because $P\succ0$, $V(x)$ is positive definite and radially unbounded.
From (1) we have $V(x_{k+1})-V(x_k)=\Delta V(x_k)<0$ for all $x_k\neq0$,
so $V$ is strictly decreasing along non‑zero trajectories.
Thus $V(x_k)$ is a monotonically decreasing sequence bounded below by
$0$, hence converges.  LaSalle’s invariance principle then implies that
the only invariant set where $\Delta V(x)=0$ is $\{0\}$; consequently
$x_k\to0$ as $k\to\infty$.  Hence the origin is asymptotically stable.
\end{proof}

The certificate returned by \texttt{certify\_l5} uses the scalar
\[
  \lambda_{\max}\bigl(K^\top P K-P\bigr)
\]
as a certificate of decrease.  The following norm bound is useful for
implementation.

\begin{lemma}[Operator‑norm bound on the Lyapunov decrease]\label{lem:lyapnorm}
Let $\|\cdot\|_2$ denote the spectral (operator) norm.
If
\[
  \|K\|_2 < \frac{1}{\sqrt{\|P\|_2\,\|P^{-1}\|_2}},
\tag{2}
\]
then $K^\top P K-P\prec0$.
\end{lemma}
\begin{proof}
From the sub‑multiplicativity of the spectral norm,
\[
  \|K^\top P K\|_2\le \|K\|_2^2\,\|P\|_2 .
\]
Using $\|P^{-1}\|_2 = 1/\lambda_{\min}(P)$ and
$\lambda_{\min}(P)\,\|P^{-1}\|_2=1$, we obtain
\[
  \lambda_{\max}(K^\top P K)\le \|K\|_2^2\,\|P\|_2 .
\]
Since $P\succ0$,
\[
  \lambda_{\max}(K^\top P K-P)\le
  \|K\|_2^2\,\|P\|_2-\lambda_{\min}(P)
  =\|P\|_2\Bigl(\|K\|_2^2-\frac{1}{\|P^{-1}\|_2}\Bigr).
\]
Thus the right‑hand side of (2) guarantees the eigenvalues of
$K^\top P K-P$ are strictly negative. ∎
\end{lemma}

\subsection{Tamper‑evidence of the HMAC‑chained audit trail}
\begin{theorem}[Tamper detection]\label{thm:tamper}
Let an audit trail consist of entries
$e_i=(\text{prime}_i,\text{cert}_i,w_i,s_i,h_i)$ with
$h_i=\mathrm{HMAC}_{k_i}(m_i)$,
$k_i=\mathtt{bytes.fromhex}(h_{i-1})$,
$m_i=\mathtt{JSON\_CANONICAL}(\text{prime}_i,\text{cert}_i,w_i,s_i)$,
and $h_0$ the fixed genesis hash.  If any field of a single entry
$e_j$ is altered, then
$\mathtt{verify\_chain}()$ returns \texttt{False}.
\end{theorem}
\begin{proof}
Assume a modification occurs at entry $j$, producing a
modified entry $\tilde e_j\neq e_j$ while all other entries remain
unchanged.  The verifier recomputes the message $\tilde m_j$ from
$\tilde e_j$ and the corresponding digest
$\tilde h_j=\mathrm{HMAC}_{\tilde k_j}(\tilde m_j)$ where
$\tilde k_j=\mathtt{bytes.fromhex}(h_{j-1})$ (unchanged because
$h_{j-1}$ depends only on earlier entries).  Since the HMAC function
is deterministic and the key $\tilde k_j$ is unchanged,
$\tilde h_j\neq h_j$ iff $\tilde m_j\neq m_j$.  Because the
modification changed at least one field of $e_j$, the canonical
JSON representation differs, hence $\tilde m_j\neq m_j$ and thus
$\tilde h_j\neq h_j$.  When the verifier reaches entry $j+1$ it uses
the stored hash $h_j$ as the key material for the next step; the
mismatch propagates and causes a comparison failure at the latest
by the verification of entry $j+1$.  Therefore
$\mathtt{verify\_chain}()=\texttt{False}$. ∎
\end{proof}

\begin{remark}
The theorem relies on the collision resistance of HMAC‑SHA256: for a
fixed key, finding two distinct messages with the same digest is
computationally infeasible.  This is the standard security assumption
for HMAC \cite{web:16}.
\end{remark}

\subsection{Operator‑norm bounds for the Level‑0 to Level‑4 certifiers}
The ACE levels L0–L4 each return a certificate with a Lipschitz
upper bound equal to the operator norm of the contraction matrix
$K$ (or a related matrix).  The following lemma guarantees that the
certified Lipschitz constant is an upper bound on the true gain.

\begin{lemma}[Spectral norm bounds the induced $2$-gain]\label{lem:opnorm}
For any matrix $M\in\mathbb{R}^{n\times n}$ and any vector $x$,
\[
  \|Mx\|_2\le\|M\|_2\,\|x\|_2 .
\]
Consequently, if a certificate reports
$\texttt{lipschitz\_upper}=\|K\|_2$, then for all inputs
\[
  \frac{\|\,Kx\,\|_2}{\|x\|_2}\le\texttt{lipschitz\_upper}.
\]
\end{lemma}
\begin{proof}
Immediate from the definition of the spectral norm as the
maximum singular value. ∎
\end{lemma}

\subsection{Scope binding and collision resistance}
The witness signature message $M$ (see Section~\ref{sec:l0-scope}) is a
concatenation of length‑prefixed fields.  The following lemma shows
that altering any scope‑binding field without detection would require
finding a collision in the underlying hash function.

\begin{lemma}[Binding via cryptographic hash]\label{lem:hashbind}
Let $H$ be a collision‑resistant hash function (e.g. SHA‑256).
If a signer produces a valid signature $\sigma=\mathrm{Sign}_{sk}(M)$
on a message $M$ that contains a field $F$ (such as \texttt{scopeHash},
\texttt{aud.class}, \texttt{nonce}, etc.), then any verifier that
reconstructs $M'$ from the signature and accepts $\sigma$ must have
$F'=F$ with overwhelming probability; otherwise a collision in $H$
would have been found.
\end{lemma}
\begin{proof}
The signature verification checks $\mathrm{Verify}_{pk}(M,\sigma)=\texttt{True}$.
Since the signature scheme is existentially unforgeable under
chosen‑message attacks, the only way to produce a valid signature on a
message $M'\neq M$ is to find a collision in the hash that the
signature scheme internally uses (or to forge the signature directly,
which is assumed infeasible).  All scope‑binding fields are fed
into $M$ via a cryptographic hash (either directly as a field or as
part of a larger hash such as \texttt{intentHash} or \texttt{scopeHash}).
Thus changing any of those fields without detection would yield a
collision in $H$, contradicting its collision resistance. ∎
\end{lemma}

These lemmas and theorems together constitute the mathematical
foundations of the claims made in the main text regarding stability,
audit‑trail integrity, and scope‑bound non‑repudiation.  They can be
included verbatim in a formal specification or security proof
artifact. 

\nocite{*}
\bibliographystyle{unsrt}  
\bibliography{references}
\end{document}